\section{Step 11 Copy the sixteen original words}\label{chacha-s11}

Copy one word at matching offsets from the input to the working output. Advance the byte offset by four until it reaches sixty-four. The destination now holds the initial working state, while the immutable input remains available for feed-forward.

\subsection{Write and inspect}
Replace the body of region \texttt{scalar\_copy}. Retain its \texttt{\# BEGIN} and \texttt{\# END} comments and the supplied labels outside them.

\Needspace{13\baselineskip}
\MARSCodeLabel{MIPS assembly}
\begin{minted}[samepage=true]{mars-idea}
    daddu $t1, $s0, $t0
    lw $t2, 0($t1)
    daddu $t1, $s1, $t0
    sw $t2, 0($t1)
    daddiu $t0, $t0, 4
    sltiu $t1, $t0, 64
    bnez $t1, scalar_copy
    nop
\end{minted}

Use recipe \texttt{s11-1} with the cumulative program through this edit. Inspect its named checkpoint before continuing to the temporary or normal return.

\subsection{Expected observation}
All sixteen output words equal the input state; the copy offset is sixty-four.

\Needspace{9\baselineskip}
\subsection{Check your understanding}
Why does the offset advance by four?

\begin{itemize}
\item \textbf{A} There are four rounds.
\item \textbf{B} Every address is four bits wide.
\item \textbf{C} Each algorithm word occupies four bytes.
\end{itemize}
